\subsection{Metric and diagnostics}
We follow the organizer's published self-evaluation notebook. Every prediction set is compared separately with each annotator view of the same physical disk. A match requires instance IoU strictly above 0.5. We sum matched IoUs and pool counts over annotator views:
\[
 \mathrm{PQ}=\frac{\sum_{(p,t)\in\mathrm{TP}}\operatorname{IoU}(p,t)}
 {|\mathrm{TP}|+0.5|\mathrm{FP}|+0.5|\mathrm{FN}|}.
\]
This is the organizer's implementation of the PQ approach \cite{kirillov2019panoptic}. Independent full-resolution fixtures verified perfect masks, misses, splits, merges, the strict threshold, and the notebook treatment of images with no predictions. We distinguish the organizer's any-overlap split and merge counts from diagnostic counts that require substantial overlap.

Thin and faint cutoffs are fixed from development annotations before confirmation. Thin means median skeleton width at most 6.32 pixels. Faint means a local flank-normalized contrast at most 2.02. Near limb means median projected spine radius at least 0.8 times the fitted apparent disk radius. Those slices overlap. The annotation procedure omits some small, faint and near-limb cases, so conclusions apply to the available annotations.

\subsection{Matched results and selection}

\begin{table}[h]\centering\small

\caption{Matched instance scoring. Each result is pooled over independent views.}

\begin{tabular}{llrrr}

\toprule

Partition & Method & PQ & FP & FN\\

\midrule

Development & Baseline & 0.3365 & 655 & 699\\

Development & Contrast & 0.3466 & 524 & 719\\

Confirmation & Baseline & 0.2931 & 472 & 480\\

Confirmation & Contrast & 0.2978 & 405 & 492\\

Udaipur & Baseline & 0.3385 & 584 & 512\\

Udaipur & Contrast & 0.3500 & 476 & 525\\

\bottomrule

\end{tabular}

\end{table}

\begin{table}[h]\centering\small

\caption{Diagnostic errors for baseline \(\rightarrow\) contrast.}

\begin{tabular}{lrrrr}

\toprule

Partition & Splits & Merges & Zero overlap & New misses\\

\midrule

Development & 73\(\to\)67 & 12\(\to\)11 & 298\(\to\)189 & 37\\

Confirmation & 87\(\to\)79 & 5\(\to\)5 & 177\(\to\)126 & 16\\

Udaipur & 70\(\to\)61 & 6\(\to\)6 & 274\(\to\)188 & 27\\

\bottomrule

\end{tabular}

\end{table}

Development: 115 disks, 186 views, 1375 annotated instances; PQ change +0.0101; paired 95\% bootstrap interval [+0.0060, +0.0148].

Confirmation: 80 disks, 127 views, 846 annotated instances; PQ change +0.0047; paired 95\% bootstrap interval [-0.0045, +0.0117].

Udaipur: 101 disks, 165 views, 1065 annotated instances; PQ change +0.0115; paired 95\% bootstrap interval [+0.0064, +0.0170].

\par\smallskip

The baseline was selected before opening confirmation. The least aggressive contrast floor, 1.8, raised development PQ but created 37 complete misses, exceeding the predeclared cap of seven. We carried this rejected setting into confirmation and the site holdout as a fixed diagnostic. Removing a component can reduce split counts without connecting a filament; this challenger never adds pixels.

\subsection{Thin, faint and near-limb instances}

\begin{table}[h]\centering\small

\caption{Confirmation match recall at instance IoU \(>0.5\).}

\begin{tabular}{lrrr}

\toprule

Slice & Truth masks & Baseline & Contrast\\

\midrule

Thin & 209 & 27.3\% & 26.8\%\\

Faint & 221 & 22.2\% & 18.1\%\\

Near limb & 107 & 57.0\% & 50.5\%\\

\bottomrule

\end{tabular}

\end{table}

On confirmation, the 0.0047 PQ gain fell below the predeclared 0.005 threshold, with 16 new complete misses against a cap of five. The contrast filter removes weak foreground responses, including examples with no annotation overlap, but it also removes genuine faint spines. The selected baseline preserves those detections (Fig.~2). Its remaining errors include isolated dark structures, disconnected portions of annotated filaments and under-coverage of faint structures. Some cases are ambiguous across annotators. Neither the foreground model nor the filter establishes physical connectivity.

The host identified PQ above 0.35 as scientifically useful in its August 20 announcement. That figure is a reference point rather than a finalist threshold. Reche and Cid \cite{reche2026} report held-out MAGFiLO crop IoU 0.70 and Dice 0.81 using a U-Net trained on MAGFiLO and SWEFil. Their crop results have different training inputs and evaluate mask overlap with a localized target. The numbers do not translate to full-disk instance PQ. Diercke et al.\ \cite{diercke2024} document seeing-related fragmentation and confusion from other dark solar structures, motivating the split and false-positive diagnostics reported here.


\subsection{Inference and reproducibility}

The selected export processed 180 supplied test images in 76.05 seconds on the RTX 4090, with 1.15 GiB peak process RAM and 1.11 GiB peak reserved GPU memory. It contains 1,293 predicted instances and 3 disks with no emitted rows. The 324,704-byte CSV fits the 10 MB form limit. The released notebook independently decodes the exact output and checks its SHA-256.

In an independent clean process, the notebook regenerated identical CSV and coverage hashes and independently decoded every mask. The validation scores above are local. No Kaggle score was available when this report was prepared.

